
Differential geometry provides the mathematical foundation for understanding and controlling nonlinear systems, a crucial capability in modern engineering and science. While linear systems offer valuable approximations, many real-world phenomena - from spacecraft attitude control to robot manipulation to biological cell regulation - are inherently nonlinear. Understanding these systems requires the elegant mathematical framework that differential geometry provides.

Consider a simple example: a satellite maintaining its orientation in space. Unlike a linear system where effects are proportional to inputs, the satellite's rotation involves complex geometric relationships between angular momentum, torque, and orientation. This nonlinearity makes precise control challenging yet essential for mission success.

The mathematical tools we'll develop in this text - smooth manifolds, tangent spaces, vector fields, and more - provide the language to describe such nonlinear systems and design controllers for them. These tools have enabled breakthroughs in:

\begin{itemize}
\item Aerospace Engineering: Precision control of aircraft and spacecraft orientation
\item Robotics: Planning smooth motions for robotic arms and autonomous vehicles
\item Biology: Modeling cellular networks and biological control systems
\item Machine Learning: Understanding the geometry of optimization landscapes
\end{itemize}

Two fundamental concepts that we'll explore in depth are nonlinear observability (can we determine a system's state from measurements?) and controllability (can we drive the system to desired states?). These properties, while straightforward for linear systems, require sophisticated geometric tools in the nonlinear case.

This text aims to build these tools systematically, starting with basic differential geometric concepts and progressing to their applications in control theory. While we'll maintain mathematical rigor, our focus will be on developing intuition and highlighting practical applications.

In the following sections, we'll review essential concepts from differential geometry: smooth manifolds, coordinate charts, tangent spaces, and flows generated by vector fields. Readers seeking additional depth may consult standard references such as Nijmeijer \& van der Schaft (1990) for control applications or Boothby (2002) for geometric foundations. Abraham \& Marsden (1978) provides an advanced treatment suitable for deeper study.


Nonlinear control is a beautiful area within mathematics, with myriad applications in fields ranging from electrical and aerospace engineering to biology and machine learning (cite). [mention nonlinear observability and controllability]

We briefly review some basic notions from differential geometry: smooth manifolds, coordinate charts, local representatives of objects (maps, vectors, etc.), vectors as differential operators, tangent spaces, and flows generated by vector fields. A more complete treatment can be found in any standard nonlinear control or differntial geometry text, such as \cite{nijmeijer_nonlinear_1990,boothby_2002}. Another excellent source is \cite{abraham_manifolds_1988}, though this is better used as a second course or a reference text. %make it feel like one is taking too many detours from the ``goal'' of getting to core results like Frobenius' theorem. 
